% preamble.tex - Praeambel des vertieften Filters Daimler Truck
% Uebernommen aus Business_Analysen/Dual_Champions, auf einen einzelnen Titel zurueckgebaut.
\usepackage[T1]{fontenc}
\usepackage[utf8]{inputenc}
\usepackage{lmodern}
\usepackage[ngerman]{babel}
\usepackage{csquotes}
\usepackage[patch={item,toc,eqnum}]{microtype}
\usepackage{xcolor}
\usepackage{amsmath}
\usepackage{amssymb}
\usepackage{booktabs}
\usepackage{longtable}
\usepackage{array}
\usepackage{tabularx}
\usepackage{siunitx}
\usepackage{graphicx}
\usepackage{tikz}
\usepackage{placeins}
\usepackage{needspace}
\usepackage[hidelinks]{hyperref}
\usepackage{url}

\areaset[current]{463.07pt}{636.6pt}
\clubpenalty=10000
\widowpenalty=10000
\displaywidowpenalty=10000
\usepackage[para]{footmisc}

\usepackage[automark]{scrlayer-scrpage}
\pagestyle{scrheadings}
\renewcommand{\sectionmark}[1]{\markright{\thesection\quad #1}}
\clearpairofpagestyles
\setkomafont{pagehead}{\normalfont\small}
\ihead{\rightmark}
\ohead{Daimler Truck}
\KOMAoptions{headsepline=true}
\cfoot*{\pagemark}
\setcounter{tocdepth}{2}

\sisetup{
  output-decimal-marker = {,}, group-separator = {.}, group-minimum-digits = 4,
  group-digits = integer, detect-weight = true, detect-family = true
}

\definecolor{quellenlink}{RGB}{20,60,120}
\newcommand{\quellensprung}[2]{\hyperref[#1]{\textcolor{quellenlink}{#2}}}
\newcommand{\kurzDaimlerA}{Daimler Truck, Annual Report 2025}
\newcommand{\kurzDaimlerB}{Daimler Truck, Annual Report 2024}
\newcommand{\kurzDaimlerC}{Daimler Truck, Annual Report 2023}
\newcommand{\kurzDaimlerD}{Daimler Truck, Annual Report 2022}
\newcommand{\Q}[2]{\footnote{\quellensprung{quelle:#1}{\csname kurz#1\endcsname}, S.~#2.}}
\newcommand{\QW}[3]{\footnote{\quellensprung{#2}{#1} (abgerufen am #3).}}

\newcommand{\Mio}[2]{#1\,Mio.\,#2}
\newcommand{\je}[2]{#1\,#2}
\newcommand{\Pz}[1]{#1\,\%}

\definecolor{vdrot}{RGB}{155,25,30}
\newcommand{\vd}[1]{{\color{vdrot}#1}}
\definecolor{bkblau}{RGB}{20,60,150}
\newcommand{\bk}[1]{{\color{bkblau}#1}}

\makeatletter
\newcommand{\tabellenkoerper}[1]{\@@input #1.tex }
\makeatother
\newcommand{\annahme}[1]{%
  \par\medskip\noindent\fbox{\parbox{0.97\linewidth}{\small\textbf{Annahme:} #1}}\par\medskip}

\newcommand{\reihentabelle}{%
\begin{table}[htbp]\centering\small
\caption{Mehrjahresreihe 2022--2025 (Mio.~EUR; Marge und ROCE in \%, EPS in EUR, Absatz in
Tsd.\ Fahrzeugen). \enquote{Ind.} = Industrial Business (ohne Financial Services). Freier
Zufluss ist der von Daimler Truck selbst ausgewiesene \enquote{Free cash flow of the
Industrial Business}, nicht selbst gerechnet (Methode S.~\pageref{sec:methode}).}\label{tab:reihe}
\begin{tabular}{@{}lrrrrrrrl@{}}\toprule
Jahr & Umsatz Ind. & EBIT ber. & Marge ber. Ind. & ROCE Ind. & FCF Ind. & EPS & Absatz & Beleg\\\midrule
\tabellenkoerper{data/reihe}
\bottomrule\end{tabular}\end{table}}

\newcommand{\kurstabelle}{%
\begin{table}[htbp]\centering\small
\caption{Absatz gegen Kurs: beide als Ver\"anderung zum Vorjahr. Schlusskurse zum
Jahresende, Yahoo Finance.}\label{tab:kurs}
\begin{tabular}{@{}lrrrr@{}}\toprule
Jahr & Absatz (Tsd.\ Fzg.) & Ver\"and.\ in \% & Kurs Jahresende (EUR) & Ver\"and.\ in \%\\\midrule
\tabellenkoerper{data/kurs}
\bottomrule\end{tabular}\end{table}}

\newcommand{\kgvtabelle}{%
\begin{table}[htbp]\centering\small
\caption{Kurs-Gewinn-Verh\"altnis zum Jahresende.}\label{tab:kgv}
\begin{tabular}{@{}lrrr@{}}\toprule
Jahr & Kurs Jahresende (EUR) & EPS (EUR) & KGV\\\midrule
\tabellenkoerper{data/kgv}
\bottomrule\end{tabular}\end{table}}

\newcommand{\ausschuettungstabelle}{%
\begin{table}[htbp]\centering\small
\caption{Verwendung des freien Zuflusses der Industrial Business (Mio.~EUR).}\label{tab:ausschuettung}
\begin{tabular}{@{}lrrrr@{}}\toprule
Jahr & Dividende & R\"uckkauf & Freier Zufluss Ind. & Anteil\\\midrule
\tabellenkoerper{data/ausschuettung}
\bottomrule\end{tabular}\end{table}}

\newcommand{\schwellentabelle}{%
\begin{table}[htbp]\centering\small
\caption{Schwellenkurs je Aktie (EUR): Kurs, zu dem der interne Zinsfuss die H\"urde von
\Pz{\Huerde} erreicht; letzte Spalte bei \Pz{\HuerdeLang}. Zufluss konstant, Endwert =
Buchwert des Eigenkapitals 2025.}\label{tab:schwelle}
\begin{tabular}{@{}lrrrrr@{}}\toprule
Szenario & Zufluss (Mio.) & 5 Jahre & 10 Jahre & 15 Jahre & 10 J., \Pz{\HuerdeLang}\\\midrule
\tabellenkoerper{data/schwelle}
\bottomrule\end{tabular}\end{table}}

\newcommand{\umkehrtabellen}{%
\begin{table}[htbp]\centering\small
\caption{Umkehrungen zum heutigen Kurs bei \Pz{\Huerde}: erforderliches Wachstum des freien
Zuflusses p.\,a.\ (oben) und erforderlicher Endwert als Vielfaches des Buchwerts
(unten).}\label{tab:umkehr}
\begin{tabular}{@{}lrrr@{}}\toprule
Szenario & 5 Jahre & 10 Jahre & 15 Jahre\\\midrule
\tabellenkoerper{data/wachstum}
\midrule
\tabellenkoerper{data/endwert}
\bottomrule\end{tabular}\end{table}}
